%=====================================================================
% PART IV OPENER
%=====================================================================
\thispagestyle{plain}
\structuralfigures{P4.}

\begin{center}
\begin{tcolorbox}[enhanced, width=0.92\textwidth, colback=purplehl!5,
  colframe=purplehl, boxrule=0pt, leftrule=5pt, arc=3pt, top=8pt, bottom=8pt]
\large\itshape\color{purplehl!70!black}
Labels are expensive. Part IV asks what can be learned without them ---
groups, directions, maps --- and then what a handful of labels can do when they
are combined with a great many unlabelled examples. The EM algorithm at its end
returns in Part V.
\end{tcolorbox}
\end{center}

\vspace{6pt}

\dsfigh{%
\begin{tikzpicture}[
  cbox/.style={rounded corners=3pt, draw=purplehl, line width=1pt, fill=purplehl!7,
               text=purplehl!70!black, font=\small\bfseries, align=center,
               minimum width=3.4cm, minimum height=1.0cm},
  arr/.style={-{Stealth[length=2.4mm]}, purplehl!60, line width=1pt}
]
  \node[cbox] (c15) at (0,0)    {Ch.\ 15\\Clustering};
  \node[cbox] (c16) at (4.6,0)  {Ch.\ 16\\PCA and SOM};
  \node[cbox] (c17) at (9.2,0)  {Ch.\ 17\\Semi-supervised \& EM};
  \node[cbox, draw=tealhl, fill=tealhl!7, text=tealhl!70!black] (c18) at (13.8,0)
       {Ch.\ 18\\Hidden Markov};
  \draw[arr] (c15) -- (c16);
  \draw[arr] (c15.south) .. controls +(0,-1.4) and +(0,-1.4) .. (c17.south);
  \draw[arr] (c16) -- (c17);
  \draw[arr, tealhl!60, dashed] (c17) -- (c18);
  \node[font=\scriptsize\itshape, text=purplehl, align=center, fill=white,
        inner sep=2pt] at (4.6,-1.05) {EM is soft $k$-means};
  \node[font=\scriptsize\itshape, text=tealhl, align=center] at (11.5,0.8)
       {Baum--Welch is EM};
\end{tikzpicture}}%
{Reading path through Part IV. $k$-means (Chapter 15) is the hard-assignment special case
of the EM algorithm in Chapter 17, which reappears in Chapter 18 to train hidden Markov
models.}{part4map}

\newpage
